\documentclass[../../../include/open-logic-section]{subfiles}

\begin{document}

\olfileid[id]{sth}{spine}{zf}
\olsection{$\Z$ dan $\ZF$: Sebuah Tonggak Penting}

Dengan Fondasi, kita mencapai satu lagi tonggak penting. Kita telah
mempertimbangkan teori $\Zminus$ dan $\ZFminus$, yang kita katakan merupakan
teori tertentu yang ``dikurangi'' sesuatu. Sesuatu itu ialah Fondasi. Jadi:

\begin{defn}
Teori $\Z$ menambahkan Fondasi kepada $\Zminus$. Jadi, aksioma-aksiomanya
ialah Ekstensionalitas, Gabungan, Pasangan, Himpunan Kuasa, Ketakhinggaan,
Fondasi, dan semua instans skema Pemisahan.

Teori $\ZF$ menambahkan Fondasi kepada $\ZFminus$. Dengan kata lain, $\ZF$
menambahkan semua instans Penggantian kepada~$\Z$.
\end{defn}

Namun, mungkin satu pertanyaan telah terlintas dalam pikiran Anda. Jika
Regularitas ekuivalen dengan Fondasi di atas $\ZFminus$, dan pembenaran
Regularitas sudah jelas, mengapa kita harus mengambil jalan memutar dan
menjadikan Fondasi, bukan Regularitas, sebagai aksioma dasar kita?

Dengan mengesampingkan alasan historis (mengenai siapa yang merumuskan apa
dan kapan), alasan dasarnya adalah bahwa Fondasi dapat disajikan tanpa
menggunakan definisi $V_\alpha$. Definisi itu bergantung pada seluruh
pekerjaan dalam \olref[recursion]{sec}: kita perlu membuktikan Rekursi
Transfinit untuk menunjukkan bahwa definisi tersebut dapat dibenarkan.
Namun, bukti Rekursi Transfinit kita menggunakan \emph{Penggantian}. Jadi,
meskipun Fondasi dan Regularitas ekuivalen modulo $\ZFminus$, keduanya tidak
ekuivalen modulo $\Zminus$.

Sesungguhnya, persoalannya lebih drastis daripada yang disiratkan oleh
pengamatan sederhana ini. Walaupun hal itu jauh melampaui cakupan buku ini,
ternyata baik $\Zminus$ maupun $\Z$ terlalu lemah untuk mendefinisikan
$V_\alpha$. Jadi, jika Anda hanya bekerja dalam $\Z$, Regularitas (sebagaimana
kita merumuskannya) bahkan tidak \emph{bermakna}. Inilah sebabnya aksioma
resmi kita ialah Fondasi, bukan Regularitas.

Mulai sekarang, kita akan bekerja dalam $\ZF$ (kecuali dinyatakan lain),
tanpa komentar lebih lanjut.

\end{document}
